%2multibyte Version: 5.50.0.2960 CodePage: 65001
%\newtheorem{definition}[definicion]{Definici�n}
%\newtheorem{example}[ejemplo]{Ejemplo}


\documentclass[10pt]{beamer}
%%%%%%%%%%%%%%%%%%%%%%%%%%%%%%%%%%%%%%%%%%%%%%%%%%%%%%%%%%%%%%%%%%%%%%%%%%%%%%%%%%%%%%%%%%%%%%%%%%%%%%%%%%%%%%%%%%%%%%%%%%%%%%%%%%%%%%%%%%%%%%%%%%%%%%%%%%%%%%%%%%%%%%%%%%%%%%%%%%%%%%%%%%%%%%%%%%%%%%%%%%%%%%%%%%%%%%%%%%%%%%%%%%%%%%%%%%%%%%%%%%%%%%%%%%%%
\usepackage{amsfonts}
\usepackage{amsmath}
\usepackage{mathpazo}
\usepackage{hyperref}
\usepackage{multimedia}
\usepackage{times}
\usepackage[T1]{fontenc}
\usepackage{multicol}

\setcounter{MaxMatrixCols}{10}
%TCIDATA{OutputFilter=LATEX.DLL}
%TCIDATA{Version=5.50.0.2960}
%TCIDATA{Codepage=65001}
%TCIDATA{<META NAME="SaveForMode" CONTENT="1">}
%TCIDATA{BibliographyScheme=Manual}
%TCIDATA{LastRevised=Sunday, March 30, 2025 23:55:44}
%TCIDATA{<META NAME="GraphicsSave" CONTENT="32">}

\newtheorem{definicion}{Definici\'{o}n}
\newtheorem{ejemplo}{Ejemplo}
\newtheorem{ejemplos}{Ejemplos}
\newenvironment{stepenumerate}{\begin{enumerate}[<+->]}{\end{enumerate}}
\newenvironment{stepitemize}{\begin{itemize}[<+->]}{\end{itemize} }
\newenvironment{stepenumeratewithalert}{\begin{enumerate}[<+-| alert@+>]}{\end{enumerate}}
\newenvironment{stepitemizewithalert}{\begin{itemize}[<+-| alert@+>]}{\end{itemize} }
\usetheme{CambridgeUS}
\usecolortheme[RGB={0,88,127}]{structure}
\usefonttheme{professionalfonts}
\input{tcilatex}
\begin{document}

\title[MAT-016]{Razonamiento L\'{o}gico Matem\'{a}tico}
\subtitle{Unidad II: \'{A}lgebra en los N\'{u}meros Reales\\
Clase 7: Conjuntos Num\'{e}ricos y Operatoria}
\author[]{Profesor Luis Tulio Gonz\'{a}lez Alcaino \\
%EndAName
Magister en Matem\'{a}tica}
\institute{Universidad Santo Tom\'{a}s\\
Departamento Ciencias B\'{a}sicas - Talca }
\date{Semestre I \ - 2023}
\maketitle

%TCIMACRO{\TeXButton{BeginFrame}{\begin{frame}}}%
%BeginExpansion
\begin{frame}%
%EndExpansion

\QTR{frametitle}{Contenidos de la clase}

\begin{enumerate}
\item Conjuntos Num\'{e}ricos

\item Operatoria
\end{enumerate}

%TCIMACRO{\TeXButton{EndFrame}{\end{frame}}}%
%BeginExpansion
\end{frame}%
%EndExpansion

%TCIMACRO{\TeXButton{BeginFrame}{\begin{frame}}}%
%BeginExpansion
\begin{frame}%
%EndExpansion

\QTR{frametitle}{Resultado de Aprendizaje}

\begin{enumerate}
\item Resolver operatoria en los distintos conjuntos num\'{e}ricos
\end{enumerate}

%TCIMACRO{\TeXButton{EndFrame}{\end{frame}}}%
%BeginExpansion
\end{frame}%
%EndExpansion

%TCIMACRO{\TeXButton{BeginFrame}{\begin{frame}}}%
%BeginExpansion
\begin{frame}%
%EndExpansion

\QTR{frametitle}{Conjuntos Num\'{e}ricos}

%TCIMACRO{\TeXButton{BeginBlok}{\begin{block}}}%
%BeginExpansion
\begin{block}%
%EndExpansion

Se iniciar\'{a} definiendo los conjuntos de n\'{u}meros que conforman a los n%
\'{u}meros reales.

\begin{enumerate}
\item \textbf{Los N\'{u}meros Naturales} $\mathbb{N}$

Tambi\'{e}n conocidos como n\'{u}meros para contar o enteros positivos.%
\begin{equation*}
\mathbb{N=}\left\{ 1,2,3,\ldots \right\}
\end{equation*}%
%TCIMACRO{\TeXButton{Pausa}{\pause}}%
%BeginExpansion
\pause%
%EndExpansion

\item \textbf{Los N\'{u}meros Enteros} $\mathbb{Z}$

Son los n\'{u}meros reales que se denotan por $\mathbb{Z}$, y se escribe%
\begin{equation*}
\mathbb{Z=}\left\{ \ldots ,-3,-2,-1,0,1,2,3,\ldots \right\}
\end{equation*}
\end{enumerate}

%TCIMACRO{\TeXButton{EndBlok}{\end{block}}}%
%BeginExpansion
\end{block}%
%EndExpansion

%TCIMACRO{\TeXButton{EndFrame}{\end{frame}}}%
%BeginExpansion
\end{frame}%
%EndExpansion

%TCIMACRO{\TeXButton{BeginFrame}{\begin{frame}}}%
%BeginExpansion
\begin{frame}%
%EndExpansion

\QTR{frametitle}{Conjuntos Num\'{e}ricos}

%TCIMACRO{\TeXButton{BeginBlok}{\begin{block}{}}}%
%BeginExpansion
\begin{block}{}%
%EndExpansion

\begin{enumerate}
\item[3.] \textbf{Los N\'{u}meros Racionales }$\mathbb{Q}$

Los n\'{u}meros racionales son los n\'{u}meros reales que se pueden expresar
como raz\'{o}n de dos n\'{u}meros enteros. Se denota el conjunto de los n%
\'{u}meros racionales por $\mathbb{Q}$, y se escribe%
\begin{eqnarray*}
\mathbb{Q} &\mathbb{=}&\left\{ \dfrac{a}{b}:\text{donde }a\text{ y }b\text{
son n\'{u}meros enteros con }b\neq 0\right\} \\
&&\left\{ \ldots ,-\frac{3}{2},\ldots ,-\frac{5}{4},\ldots ,-1,\ldots ,-%
\frac{1}{2},\ldots ,0,\ldots ,\frac{1}{2}\ldots ,\frac{8}{5},\ldots \right\}
\end{eqnarray*}

Por otro lado, su desarrollo decimal es finito o infinito peri\'{o}dico o
semiperi\'{o}dico

\textbf{Ejemplos:}
\end{enumerate}

$\dfrac{1}{4}=\allowbreak 0.25\qquad \qquad \dfrac{5}{3}$ : $1.\allowbreak
666\ldots =1.\overline{\allowbreak 6}\qquad \qquad \dfrac{19}{12}%
=1.58333\ldots =1.58\overline{3}$

%TCIMACRO{\TeXButton{EndBlok}{\end{block}}}%
%BeginExpansion
\end{block}%
%EndExpansion

%TCIMACRO{\TeXButton{EndFrame}{\end{frame}}}%
%BeginExpansion
\end{frame}%
%EndExpansion

%TCIMACRO{\TeXButton{BeginFrame}{\begin{frame}}}%
%BeginExpansion
\begin{frame}%
%EndExpansion

\QTR{frametitle}{Conjuntos Num\'{e}ricos}

%TCIMACRO{\TeXButton{BeginBlok}{\begin{block}{}}}%
%BeginExpansion
\begin{block}{}%
%EndExpansion

\begin{enumerate}
\item[4.] \textbf{Los N\'{u}meros Irracionales }

Son los n\'{u}meros que no pueden expresarse como un cociente de enteros y
su desarrollo decimal es infinito no peri\'{o}dico.

\begin{eqnarray*}
\mathbb{I} &\mathbb{=}&\left\{ x:\text{donde }x\text{ no es un n\'{u}meros
racional}\right\} \\
&&\left\{ \ldots ,-\pi ,\ldots ,-\sqrt{3},\ldots ,\sqrt{2},\ldots ,e,\ldots
\right\}
\end{eqnarray*}

\textbf{Ejemplos:}

$\sqrt{2}=\allowbreak 1.414213562\ldots \qquad \pi =3.141592654\ldots \qquad
e=2.7182818228\ldots $
\end{enumerate}

%TCIMACRO{\TeXButton{EndBlok}{\end{block}}}%
%BeginExpansion
\end{block}%
%EndExpansion

%TCIMACRO{\TeXButton{EndFrame}{\end{frame}}}%
%BeginExpansion
\end{frame}%
%EndExpansion

%TCIMACRO{\TeXButton{BeginFrame}{\begin{frame}}}%
%BeginExpansion
\begin{frame}%
%EndExpansion

\QTR{frametitle}{Conjuntos Num\'{e}ricos}

%TCIMACRO{\TeXButton{BeginBlok}{\begin{block}{}}}%
%BeginExpansion
\begin{block}{}%
%EndExpansion

\begin{enumerate}
\item[5.] \textbf{Los N\'{u}meros Reales} $\mathbb{R}$

La uni\'{o}n del conjunto de los n\'{u}meros racionales y de los n\'{u}meros
irracionales se conoce como el conjunto de los n\'{u}meros reales, es decir, 
$\mathbb{R=Q\cup I}$
\end{enumerate}

Una manera bastante pr\'{a}ctica es representar el conjunto de n\'{u}meros
reales es en una recta a la que se denomina recta num\'{e}rica como la que
aparece en la siguiente figura

\FRAME{ftbpF}{3.7429in}{0.3649in}{0pt}{}{}{Figure}{\special{language
"Scientific Word";type "GRAPHIC";maintain-aspect-ratio TRUE;display
"USEDEF";valid_file "T";width 3.7429in;height 0.3649in;depth
0pt;original-width 8.412in;original-height 0.7948in;cropleft "0";croptop
"1";cropright "1";cropbottom "0";tempfilename
'RSJ20W03.bmp';tempfile-properties "XPR";}}

%TCIMACRO{\TeXButton{EndBlok}{\end{block}}}%
%BeginExpansion
\end{block}%
%EndExpansion

%TCIMACRO{\TeXButton{EndFrame}{\end{frame}}}%
%BeginExpansion
\end{frame}%
%EndExpansion

%TCIMACRO{\TeXButton{BeginFrame}{\begin{frame}}}%
%BeginExpansion
\begin{frame}%
%EndExpansion

\QTR{frametitle}{Conjuntos Num\'{e}ricos}

%TCIMACRO{\TeXButton{BeginBlok}{\begin{block}{}}}%
%BeginExpansion
\begin{block}{}%
%EndExpansion

\textbf{Recta num\'{e}rica}

\FRAME{ftbpF}{3.7429in}{0.3649in}{0pt}{}{}{Figure}{\special{language
"Scientific Word";type "GRAPHIC";maintain-aspect-ratio TRUE;display
"USEDEF";valid_file "T";width 3.7429in;height 0.3649in;depth
0pt;original-width 8.412in;original-height 0.7948in;cropleft "0";croptop
"1";cropright "1";cropbottom "0";tempfilename
'RSJ1NG02.bmp';tempfile-properties "XPR";}}

Mediante las ideas de igualdad y desigualdad pueden compararse 2 n\'{u}meros
reales cualquiera. Se utilizan s\'{\i}mbolos para representar estas ideas:

%TCIMACRO{\TeXButton{EndBlok}{\end{block}}}%
%BeginExpansion
\end{block}%
%EndExpansion

%TCIMACRO{\TeXButton{Pausa}{\pause}}%
%BeginExpansion
\pause%
%EndExpansion
\textbf{Ejemplos}

\begin{itemize}
\item Para decir que 7 es menor que 12 , escribimos $7<12$%
%TCIMACRO{\TeXButton{Pausa}{\pause}}%
%BeginExpansion
\pause%
%EndExpansion

\item Tambi\'{e}n para decir que $-3$ es mayor que $-10$, escribimos $-3>-10$%
%TCIMACRO{\TeXButton{Pausa}{\pause}}%
%BeginExpansion
\pause%
%EndExpansion

\item Podemos tener claro el significado de los s\'{\i}mbolos menor que $<$
\ y mayor que $>$ , si recordamos que \'{e}stos siempre apuntan hacia el n%
\'{u}mero m\'{a}s peque\~{n}o: $-5<1,$as\'{\i} como $15>6$%
%TCIMACRO{\TeXButton{Pausa}{\pause}}%
%BeginExpansion
\pause%
%EndExpansion

\item Hay otros dos s\'{\i}mbolos $\leq $ y $\geq $ , que tambi\'{e}n
representan la idea de desigualdad. El s\'{\i}mbolo $\leq $ significa
\textquotedblleft es menor o igual que\textquotedblright , por lo que $%
10\leq 15$ significa que 10 es menor o igual a 15.
\end{itemize}

%TCIMACRO{\TeXButton{EndFrame}{\end{frame}}}%
%BeginExpansion
\end{frame}%
%EndExpansion

%TCIMACRO{\TeXButton{BeginFrame}{\begin{frame}}}%
%BeginExpansion
\begin{frame}%
%EndExpansion

\QTR{frametitle}{Conjuntos Num\'{e}ricos}

%TCIMACRO{\TeXButton{BeginBlok}{\begin{block}{Ejercicio}}}%
%BeginExpansion
\begin{block}{Ejercicio}%
%EndExpansion

Complete el siguiente cuadro: Marca con \textquotedblleft $\surd $%
\textquotedblright\ el (o los) conjunto(s) num\'{e}rico(s) al(los) cual(es)
pertenece cada n\'{u}mero.

\begin{equation*}
\begin{tabular}{|c|c|c|c|c|c|}
\hline
N\'{u}mero & $\mathbb{N}$ & $\mathbb{Z}$ & $\mathbb{Q}$ & $\mathbb{I}$ & $%
\mathbb{R}$ \\ \hline
$-5$ &  &  &  &  &  \\ \hline
$\sqrt{3}$ &  &  &  &  &  \\ \hline
$\frac{5}{8}$ &  &  &  &  &  \\ \hline
$\frac{\pi }{2}$ &  &  &  &  &  \\ \hline
$1,2$ &  &  &  &  &  \\ \hline
\end{tabular}%
\end{equation*}

%TCIMACRO{\TeXButton{EndBlok}{\end{block}}}%
%BeginExpansion
\end{block}%
%EndExpansion

%TCIMACRO{\TeXButton{EndFrame}{\end{frame}}}%
%BeginExpansion
\end{frame}%
%EndExpansion

%TCIMACRO{\TeXButton{BeginFrame}{\begin{frame}}}%
%BeginExpansion
\begin{frame}%
%EndExpansion

\QTR{frametitle}{Conjuntos Num\'{e}ricos}

%TCIMACRO{\TeXButton{BeginBlok}{\begin{block}{Adici\'{o}n}}}%
%BeginExpansion
\begin{block}{Adici\'{o}n}%
%EndExpansion

Para sumar N\'{u}meros Enteros, debemos tener en cuenta que:

\begin{enumerate}
\item \textquotedblleft Si son del mismo signo (positivos o negativos), se
suman los n\'{u}meros (sin considerar el signo) y se conseva el signo de los
n\'{u}meros (positivo o negativo)\textquotedblright

\textbf{Ejemplos:}

\begin{enumerate}
\item $2+5=7$

\item $-4+(-7)=-11$%
%TCIMACRO{\TeXButton{Pausa}{\pause}}%
%BeginExpansion
\pause%
%EndExpansion
\end{enumerate}

\item \textquotedblleft Si son de distinto signo (positivos o negativos o
viceversa), se restan los n\'{u}meros y se conseva el signo (positivo o
negativo) del n\'{u}mero mayor sin considerar el signo\textquotedblright

\textbf{Ejemplos:}

\begin{enumerate}
\item $8+(-3)=5$

\item $-10+9=-1$
\end{enumerate}
\end{enumerate}

%TCIMACRO{\TeXButton{EndBlok}{\end{block}}}%
%BeginExpansion
\end{block}%
%EndExpansion

%TCIMACRO{\TeXButton{EndFrame}{\end{frame}}}%
%BeginExpansion
\end{frame}%
%EndExpansion

%TCIMACRO{\TeXButton{BeginFrame}{\begin{frame}}}%
%BeginExpansion
\begin{frame}%
%EndExpansion

\QTR{frametitle}{Conjuntos Num\'{e}ricos}

%TCIMACRO{\TeXButton{BeginBlok}{\begin{block}{Observaciones}}}%
%BeginExpansion
\begin{block}{Observaciones}%
%EndExpansion

\begin{enumerate}
\item La expresi\'{o}n \textquotedblleft $a+(-b)$\textquotedblright\ se
escribe como \textquotedblleft $a-b$\textquotedblright

\textbf{Ejemplo.} $6+(-8)=6-8=-2$%
%TCIMACRO{\TeXButton{Pausa}{\pause}}%
%BeginExpansion
\pause%
%EndExpansion

\item La expresi\'{o}n \textquotedblleft $a-(-b)$\textquotedblright\ se
escribe como \textquotedblleft $a+b$\textquotedblright

\textbf{Ejemplo. }\ $12-(-4)=12+4=16$%
%TCIMACRO{\TeXButton{Pausa}{\pause}}%
%BeginExpansion
\pause%
%EndExpansion

\item La expresi\'{o}n \textquotedblleft $a+b$\textquotedblright\ es
equivalente con \textquotedblleft $b+a$\textquotedblright $,$ es decir $%
a+b=b+a$. (propiedad conmutativa).

\textbf{Ejemplo.} $(-25)+10=10+(-25)=10-25=-15$%
%TCIMACRO{\TeXButton{Pausa}{\pause}}%
%BeginExpansion
\pause%
%EndExpansion

\item La expresi\'{o}n \textquotedblleft $a+(b+c)$\textquotedblright\ es
equivalente con \textquotedblleft $(a+b)+c$\textquotedblright $,$ es decir $%
a+(b+c)=(a+b)+c$. (propiedad asociativa).

\textbf{Ejemplo.}

$((-12)+15)+(-3)=\allowbreak (-12)+(15+(-3))=(-12)+(15-3)=-12+12=0$
\end{enumerate}

%TCIMACRO{\TeXButton{EndBlok}{\end{block}}}%
%BeginExpansion
\end{block}%
%EndExpansion

%TCIMACRO{\TeXButton{EndFrame}{\end{frame}}}%
%BeginExpansion
\end{frame}%
%EndExpansion

%TCIMACRO{\TeXButton{BeginFrame}{\begin{frame}}}%
%BeginExpansion
\begin{frame}%
%EndExpansion

\QTR{frametitle}{Conjuntos Num\'{e}ricos}

%TCIMACRO{\TeXButton{BeginBlok}{\begin{block}{Observaciones}}}%
%BeginExpansion
\begin{block}{Observaciones}%
%EndExpansion

\begin{enumerate}
\item $a+0=a$.

\textbf{Ejemplo. }$7+0=7$%
%TCIMACRO{\TeXButton{Pausa}{\pause}}%
%BeginExpansion
\pause%
%EndExpansion

\item $a+(-a)=(-a)+a=0$.

\textbf{Ejemplo. }$12+(-12)=0$%
%TCIMACRO{\TeXButton{Pausa}{\pause}}%
%BeginExpansion
\pause%
%EndExpansion

\item $-(-a)=a$.%
%TCIMACRO{\TeXButton{Pausa}{\pause}}%
%BeginExpansion
\pause%
%EndExpansion

\textbf{Ejemplo. }$-(-9)=9$

\item $-(a+b)=(-a)+(-b)$%
%TCIMACRO{\TeXButton{Pausa}{\pause}}%
%BeginExpansion
\pause%
%EndExpansion

\item \textbf{Ejemplo. }$-(4+8)=(-4)+(-8)=\allowbreak -12$
\end{enumerate}

%TCIMACRO{\TeXButton{EndBlok}{\end{block}}}%
%BeginExpansion
\end{block}%
%EndExpansion

%TCIMACRO{\TeXButton{EndFrame}{\end{frame}}}%
%BeginExpansion
\end{frame}%
%EndExpansion

%TCIMACRO{\TeXButton{BeginFrame}{\begin{frame}}}%
%BeginExpansion
\begin{frame}%
%EndExpansion

\QTR{frametitle}{Conjuntos Num\'{e}ricos}

%TCIMACRO{\TeXButton{BeginBlok}{\begin{block}{}}}%
%BeginExpansion
\begin{block}{}%
%EndExpansion

\textbf{Ejercicios}

\begin{enumerate}
\item Determine el valor de las siguientes sumas:

\begin{tabular}{llll}
a) $(-3)+(-8)=$%
%TCIMACRO{\TeXButton{Pausa}{\pause} }%
%BeginExpansion
\pause
%EndExpansion
&  & b) $(-5+3)+2=$%
%TCIMACRO{\TeXButton{Pausa}{\pause} }%
%BeginExpansion
\pause
%EndExpansion
&  \\ 
c) $-10+5=$%
%TCIMACRO{\TeXButton{Pausa}{\pause} }%
%BeginExpansion
\pause
%EndExpansion
&  & d) $0+(-5+2)=$%
%TCIMACRO{\TeXButton{Pausa}{\pause} }%
%BeginExpansion
\pause
%EndExpansion
&  \\ 
e) $(-32)+48+(-10)=$%
%TCIMACRO{\TeXButton{Pausa}{\pause} }%
%BeginExpansion
\pause
%EndExpansion
&  & f) $-(-14)+(-18)=$%
%TCIMACRO{\TeXButton{Pausa}{\pause} }%
%BeginExpansion
\pause
%EndExpansion
&  \\ 
g) $-(3+(-7))-7-3=$%
%TCIMACRO{\TeXButton{Pausa}{\pause} }%
%BeginExpansion
\pause
%EndExpansion
&  & h) $-(-8+5)+(-8)=$ & 
\end{tabular}

\item Completar los espacios en blanco en cada una de las siguientes
expresiones de modo que resulte una expresi\'{o}n verdadera.

\begin{enumerate}
\item $3+(-3)+6=4+\ \ \ \ \underline{\text{ \ \ \ \ \ \ \ \ \ \ \ }}\ \ \ \
\ \ \ \ $%
%TCIMACRO{\TeXButton{Pausa}{\pause}}%
%BeginExpansion
\pause%
%EndExpansion

\item $(-5)+(-(-5))+(-(-(-5)))=\underline{\ \ \ \ \ \ \ \ \ \ \ }$%
%TCIMACRO{\TeXButton{Pausa}{\pause}}%
%BeginExpansion
\pause%
%EndExpansion

\item $9=9+\underline{\ \ \ \ \ \ \ \ \ \ \ }$%
%TCIMACRO{\TeXButton{Pausa}{\pause}}%
%BeginExpansion
\pause%
%EndExpansion

\item $-8+6+\underline{\ \ \ \ \ \ \ \ \ \ \ \ \ }=-8$
\end{enumerate}
\end{enumerate}

%TCIMACRO{\TeXButton{EndBlok}{\end{block}}}%
%BeginExpansion
\end{block}%
%EndExpansion

%TCIMACRO{\TeXButton{EndFrame}{\end{frame}}}%
%BeginExpansion
\end{frame}%
%EndExpansion

%TCIMACRO{\TeXButton{BeginFrame}{\begin{frame}}}%
%BeginExpansion
\begin{frame}%
%EndExpansion

\QTR{frametitle}{Conjuntos Num\'{e}ricos}

%TCIMACRO{\TeXButton{BeginBlok}{\begin{block}{Multiplicaci\'{o}n}}}%
%BeginExpansion
\begin{block}{Multiplicaci\'{o}n}%
%EndExpansion

Para multiplicar n\'{u}meros enteros, debemos tener en cuenta la
\textquotedblleft ley de los signos\textquotedblright , que dice lo
siguiente:%
%TCIMACRO{\TeXButton{Pausa}{\pause}}%
%BeginExpansion
\pause%
%EndExpansion

\begin{itemize}
\item Al multiplicar dos n\'{u}meros enteros de igual signo, el resultado es
un n\'{u}mero positivo $(+)$%
%TCIMACRO{\TeXButton{Pausa}{\pause}}%
%BeginExpansion
\pause%
%EndExpansion

\item Al multiplicar dos n\'{u}meros enteros de distinto signo, el resultado
es un n\'{u}mero negativo $(-)$%
%TCIMACRO{\TeXButton{Pausa}{\pause}}%
%BeginExpansion
\pause%
%EndExpansion
\end{itemize}

Esto se puede resumir en la siguiente tabla:%
%TCIMACRO{\TeXButton{Pausa}{\pause}}%
%BeginExpansion
\pause%
%EndExpansion

\begin{eqnarray*}
&&\text{{\large Ley de los Signos}} \\
&&%
\begin{tabular}{ccc}
$+\cdot +$ & $=$ & $+$ \\ 
$-\cdot -$ & $=$ & $+$ \\ 
$+\cdot -$ & $=$ & $-$ \\ 
$-\cdot +$ & $=$ & $-$%
\end{tabular}%
\end{eqnarray*}

%TCIMACRO{\TeXButton{EndBlok}{\end{block}}}%
%BeginExpansion
\end{block}%
%EndExpansion

%TCIMACRO{\TeXButton{EndFrame}{\end{frame}}}%
%BeginExpansion
\end{frame}%
%EndExpansion

%TCIMACRO{\TeXButton{BeginFrame}{\begin{frame}}}%
%BeginExpansion
\begin{frame}%
%EndExpansion

\QTR{frametitle}{Conjuntos Num\'{e}ricos}

%TCIMACRO{\TeXButton{BeginBlok}{\begin{block}{Observaciones}}}%
%BeginExpansion
\begin{block}{Observaciones}%
%EndExpansion
\textbf{\ }

\begin{enumerate}
\item $a\cdot b=b\cdot a$ (conmutatividad)

\textbf{Ejemplo. }$\mathbf{(-2)\cdot 7=7\cdot (-2)}=\allowbreak -14$%
%TCIMACRO{\TeXButton{Pausa}{\pause}}%
%BeginExpansion
\pause%
%EndExpansion

\item $a\cdot \left( b\cdot c\right) =\left( a\cdot b\right) \cdot c$
(asociatividad)

\textbf{Ejemplo. }$4\cdot \left( 12\cdot (-3\right) )=(4\cdot 12\cdot (-3)$%
%TCIMACRO{\TeXButton{Pausa}{\pause}}%
%BeginExpansion
\pause%
%EndExpansion

\item $a\cdot 1=a$ (neutro multiplicativo). \textbf{Ejemplo }$83\cdot 1=83$%
%TCIMACRO{\TeXButton{Pausa}{\pause}}%
%BeginExpansion
\pause%
%EndExpansion

\item $a\cdot (b+c)=a\cdot b+a\cdot c$ (distributividad)

\textbf{Ejemplo. }$(-4)\cdot (5+7)=(-4)\cdot 5+(-4)\cdot 7=\allowbreak -48$%
%TCIMACRO{\TeXButton{Pausa}{\pause}}%
%BeginExpansion
\pause%
%EndExpansion

\item $a\cdot 0=0.$ \textbf{Ejemplo }$(-120)\cdot 0=\allowbreak 0$%
%TCIMACRO{\TeXButton{Pausa}{\pause}}%
%BeginExpansion
\pause%
%EndExpansion

\item $a\cdot b=0$, entonces $a=0$ o $b=0$, o bien ambos son iguales a cero

\textbf{Ejemplo.}

$(x-1)\cdot (x+5)=0\Longrightarrow (x-1)=0$ \ o bien \ \textbf{\ }$(x+5)=0$

de donde $x=1$ \ o bien \ \ \ $x=-5.$
\end{enumerate}

%TCIMACRO{\TeXButton{EndBlok}{\end{block}}}%
%BeginExpansion
\end{block}%
%EndExpansion

%TCIMACRO{\TeXButton{EndFrame}{\end{frame}}}%
%BeginExpansion
\end{frame}%
%EndExpansion

%TCIMACRO{\TeXButton{BeginFrame}{\begin{frame}}}%
%BeginExpansion
\begin{frame}%
%EndExpansion

\QTR{frametitle}{Conjuntos Num\'{e}ricos}

%TCIMACRO{\TeXButton{BeginBlok}{\begin{block}{Ejercicios}}}%
%BeginExpansion
\begin{block}{Ejercicios}%
%EndExpansion

Multiplicar:

\begin{tabular}{lll}
a) $(-5)\cdot 0=$%
%TCIMACRO{\TeXButton{Pausa}{\pause} }%
%BeginExpansion
\pause
%EndExpansion
&  & f) $(-3)\cdot (-7)=$%
%TCIMACRO{\TeXButton{Pausa}{\pause} }%
%BeginExpansion
\pause
%EndExpansion
\\ 
b) $((-10)\cdot 4)\cdot 5=$%
%TCIMACRO{\TeXButton{Pausa}{\pause} }%
%BeginExpansion
\pause
%EndExpansion
&  & g) $3+2\cdot 5=$%
%TCIMACRO{\TeXButton{Pausa}{\pause} }%
%BeginExpansion
\pause
%EndExpansion
\\ 
c) $((-20)+40)\cdot (-2)=$%
%TCIMACRO{\TeXButton{Pausa}{\pause} }%
%BeginExpansion
\pause
%EndExpansion
&  & h) $(-2)\cdot (-(-14)+(-5))=$%
%TCIMACRO{\TeXButton{Pausa}{\pause} }%
%BeginExpansion
\pause
%EndExpansion
\\ 
d) $4\cdot (-5)=$%
%TCIMACRO{\TeXButton{Pausa}{\pause} }%
%BeginExpansion
\pause
%EndExpansion
&  & i) $(-7+3)\cdot (-3)=$%
%TCIMACRO{\TeXButton{Pausa}{\pause} }%
%BeginExpansion
\pause
%EndExpansion
\\ 
e) $20\cdot ((-8)-6)=$%
%TCIMACRO{\TeXButton{Pausa}{\pause} }%
%BeginExpansion
\pause
%EndExpansion
&  & j) $5\cdot 4-5+6\cdot 3-2+7\cdot 3=$%
\end{tabular}

%TCIMACRO{\TeXButton{EndBlok}{\end{block}}}%
%BeginExpansion
\end{block}%
%EndExpansion

%TCIMACRO{\TeXButton{EndFrame}{\end{frame}}}%
%BeginExpansion
\end{frame}%
%EndExpansion

%TCIMACRO{\TeXButton{BeginFrame}{\begin{frame}}}%
%BeginExpansion
\begin{frame}%
%EndExpansion

\QTR{frametitle}{Conjuntos Num\'{e}ricos}

%TCIMACRO{%
%\TeXButton{BeginBlok}{\begin{block}{Jerarqu�a o Prioridad de las Operaciones}}}%
%BeginExpansion
\begin{block}{Jerarqu�a o Prioridad de las Operaciones}%
%EndExpansion

Cuando se eval\'{u}a una expresi\'{o}n que contiene m\'{a}s de un s\'{\i}%
mbolo de operaci\'{o}n, se procede en el orden que describiremos a continuaci%
\'{o}n.%
%TCIMACRO{\TeXButton{Pausa}{\pause}}%
%BeginExpansion
\pause%
%EndExpansion

Diremos que los procedimientos u operaciones que aparecen primero son los
que tienen mayor jerarqu\'{\i}a o prioridad, que aquellos que aparecen despu%
\'{e}s.%
%TCIMACRO{\TeXButton{Pausa}{\pause}}%
%BeginExpansion
\pause%
%EndExpansion

\begin{enumerate}
\item Efectuar las operaciones dentro de los par\'{e}ntesis.%
%TCIMACRO{\TeXButton{Pausa}{\pause}}%
%BeginExpansion
\pause%
%EndExpansion

\item Realizar las multiplicaciones y divisiones.%
%TCIMACRO{\TeXButton{Pausa}{\pause}}%
%BeginExpansion
\pause%
%EndExpansion

\item Realizar las sumas y restas.%
%TCIMACRO{\TeXButton{Pausa}{\pause}}%
%BeginExpansion
\pause%
%EndExpansion

\item Cuando haya dos o m\'{a}s operaciones de igual jerarqu\'{\i}a, se
procede de izquierda a derecha.%
%TCIMACRO{\TeXButton{Pausa}{\pause}}%
%BeginExpansion
\pause%
%EndExpansion
\end{enumerate}

\textbf{Ejemplos}

\begin{enumerate}
\item $2+[3-(5-2\cdot 3)\cdot 4]=%
%TCIMACRO{\TeXButton{Pausa}{\pause}}%
%BeginExpansion
\pause%
%EndExpansion
\allowbreak 9$%
%TCIMACRO{\TeXButton{Pausa}{\pause}}%
%BeginExpansion
\pause%
%EndExpansion

\item $(5\cdot 3-2)-[4\cdot (5\cdot 3)-1]=%
%TCIMACRO{\TeXButton{Pausa}{\pause}}%
%BeginExpansion
\pause%
%EndExpansion
\allowbreak -46=$%
%TCIMACRO{\TeXButton{Pausa}{\pause}}%
%BeginExpansion
\pause%
%EndExpansion
\end{enumerate}

%TCIMACRO{\TeXButton{EndBlok}{\end{block}}}%
%BeginExpansion
\end{block}%
%EndExpansion

%TCIMACRO{\TeXButton{EndFrame}{\end{frame}}}%
%BeginExpansion
\end{frame}%
%EndExpansion

%TCIMACRO{\TeXButton{BeginFrame}{\begin{frame}}}%
%BeginExpansion
\begin{frame}%
%EndExpansion

\QTR{frametitle}{Operatoria en los N\'{u}meros Racionales}

%TCIMACRO{\TeXButton{BeginBlok}{\begin{block}{Adici\'{o}n}}}%
%BeginExpansion
\begin{block}{Adici\'{o}n}%
%EndExpansion

Para sumar n\'{u}meros racionales debemos tener en cuenta que:%
%TCIMACRO{\TeXButton{Pausa}{\pause}}%
%BeginExpansion
\pause%
%EndExpansion

\begin{enumerate}
\item Si tienen denominadores com\'{u}n, se suman los numeradores y se
conserva el denominador:%
\begin{equation*}
\frac{a}{b}+\frac{c}{b}=\frac{a+c}{b}\text{ , \ }b\neq 0
\end{equation*}%
%TCIMACRO{\TeXButton{Pausa}{\pause}}%
%BeginExpansion
\pause%
%EndExpansion

\item Si tienen distinto denominador:%
\begin{equation*}
\frac{a}{b}+\frac{c}{d}=\frac{a\cdot d+b\cdot c}{b\cdot d}\text{ , \ }b\neq 0%
\text{ y \ }d\neq 0
\end{equation*}%
%TCIMACRO{\TeXButton{Pausa}{\pause}}%
%BeginExpansion
\pause%
%EndExpansion

\item M\'{\i}nimo com\'{u}n denominador:

$\dfrac{3}{5}-\dfrac{8}{3}+\dfrac{6}{12}=\dfrac{12\cdot 3-20\cdot 8+5\cdot 6%
}{60}=%
%TCIMACRO{\TeXButton{Pausa}{\pause}}%
%BeginExpansion
\pause%
%EndExpansion
=\dfrac{36-160+30}{60}=\allowbreak \dfrac{-94}{60}=\allowbreak -\dfrac{47}{30%
}$
\end{enumerate}

%TCIMACRO{\TeXButton{EndBlok}{\end{block}}}%
%BeginExpansion
\end{block}%
%EndExpansion

%TCIMACRO{\TeXButton{EndFrame}{\end{frame}}}%
%BeginExpansion
\end{frame}%
%EndExpansion

%TCIMACRO{\TeXButton{BeginFrame}{\begin{frame}}}%
%BeginExpansion
\begin{frame}%
%EndExpansion

\QTR{frametitle}{Conjuntos Num\'{e}ricos}

%TCIMACRO{\TeXButton{BeginBlok}{\begin{block}{Ejemplos}}}%
%BeginExpansion
\begin{block}{Ejemplos}%
%EndExpansion

\begin{enumerate}
\item $\dfrac{1}{2}-\dfrac{5}{2}+\dfrac{3}{2}=\dfrac{1-5+3}{2}=%
%TCIMACRO{\TeXButton{Pausa}{\pause}}%
%BeginExpansion
\pause%
%EndExpansion
-\dfrac{1}{2}$%
%TCIMACRO{\TeXButton{Pausa}{\pause}}%
%BeginExpansion
\pause%
%EndExpansion

\item $\dfrac{1}{3}-\dfrac{2}{5}=\dfrac{1\cdot 5-3\cdot 2}{3\cdot 5}=\dfrac{%
5-6}{15}=%
%TCIMACRO{\TeXButton{Pausa}{\pause}}%
%BeginExpansion
\pause%
%EndExpansion
\allowbreak -\dfrac{1}{15}$%
%TCIMACRO{\TeXButton{Pausa}{\pause}}%
%BeginExpansion
\pause%
%EndExpansion

\item $\dfrac{4}{5}+\dfrac{8}{3}-\dfrac{7}{18}+\dfrac{9}{2}=\dfrac{18\cdot
4+30\cdot 8-5\cdot 7+45\cdot 9}{90}\allowbreak =\dfrac{682}{90}=\dfrac{341}{%
45}$%
%TCIMACRO{\TeXButton{Pausa}{\pause}}%
%BeginExpansion
\pause%
%EndExpansion

\item $\dfrac{1}{4}+\dfrac{7}{5}-\dfrac{3}{8}=%
%TCIMACRO{\TeXButton{Pausa}{\pause}}%
%BeginExpansion
\pause%
%EndExpansion
\dfrac{51}{40}$
\end{enumerate}

%TCIMACRO{\TeXButton{EndBlok}{\end{block}}}%
%BeginExpansion
\end{block}%
%EndExpansion

%TCIMACRO{\TeXButton{EndFrame}{\end{frame}}}%
%BeginExpansion
\end{frame}%
%EndExpansion

%TCIMACRO{\TeXButton{BeginFrame}{\begin{frame}}}%
%BeginExpansion
\begin{frame}%
%EndExpansion

\QTR{frametitle}{Conjuntos Num\'{e}ricos}

%TCIMACRO{\TeXButton{BeginBlok}{\begin{block}{Observaciones}}}%
%BeginExpansion
\begin{block}{Observaciones}%
%EndExpansion

\begin{enumerate}
\item $\dfrac{-a}{b}=\dfrac{a}{-b}=-\dfrac{a}{b}$%
%TCIMACRO{\TeXButton{Pausa}{\pause}}%
%BeginExpansion
\pause%
%EndExpansion

\item $\dfrac{-a}{-b}=\dfrac{a}{b}$%
%TCIMACRO{\TeXButton{Pausa}{\pause}}%
%BeginExpansion
\pause%
%EndExpansion

\item $-\left( \dfrac{a}{b}\right) =-\dfrac{a}{b}$%
%TCIMACRO{\TeXButton{Pausa}{\pause}}%
%BeginExpansion
\pause%
%EndExpansion

\item $\dfrac{a}{1}=a$%
%TCIMACRO{\TeXButton{Pausa}{\pause}}%
%BeginExpansion
\pause%
%EndExpansion

\item $\dfrac{a}{a}=1$%
%TCIMACRO{\TeXButton{Pausa}{\pause}}%
%BeginExpansion
\pause%
%EndExpansion

\item $\dfrac{0}{a}=0$%
%TCIMACRO{\TeXButton{Pausa}{\pause}}%
%BeginExpansion
\pause%
%EndExpansion

\item $\dfrac{a}{0}$ NO EXISTE
\end{enumerate}

%TCIMACRO{\TeXButton{EndBlok}{\end{block}}}%
%BeginExpansion
\end{block}%
%EndExpansion

%TCIMACRO{\TeXButton{EndFrame}{\end{frame}}}%
%BeginExpansion
\end{frame}%
%EndExpansion

%TCIMACRO{\TeXButton{BeginFrame}{\begin{frame}}}%
%BeginExpansion
\begin{frame}%
%EndExpansion

\QTR{frametitle}{Conjuntos Num\'{e}ricos}

%TCIMACRO{\TeXButton{BeginBlok}{\begin{block}{Multiplicaci\'{o}n}}}%
%BeginExpansion
\begin{block}{Multiplicaci\'{o}n}%
%EndExpansion

Para multiplicar dos n\'{u}meros racionales se multiplica numerador por
numerador y denominador con denominador, esto es:%
%TCIMACRO{\TeXButton{Pausa}{\pause}}%
%BeginExpansion
\pause%
%EndExpansion

\begin{equation*}
\frac{a}{b}\cdot \frac{c}{d}=\frac{a\cdot c}{b\cdot d}\text{ , \ }b\neq 0%
\text{, }d\neq 0
\end{equation*}%
%TCIMACRO{\TeXButton{Pausa}{\pause}}%
%BeginExpansion
\pause%
%EndExpansion

\textbf{Ejemplos.}

\begin{enumerate}
\item $\dfrac{1}{5}\cdot \dfrac{(-3)}{8}=\dfrac{1\cdot (-3)}{5\cdot 8}=%
%TCIMACRO{\TeXButton{Pausa}{\pause}}%
%BeginExpansion
\pause%
%EndExpansion
\dfrac{-3}{40}=$%
%TCIMACRO{\TeXButton{Pausa}{\pause}}%
%BeginExpansion
\pause%
%EndExpansion
$-\dfrac{3}{40}$

\item $\dfrac{3}{\left( -8\right) }\cdot \dfrac{5}{7}=\allowbreak =\dfrac{%
3\cdot 5}{\left( -8\right) \cdot 7}=%
%TCIMACRO{\TeXButton{Pausa}{\pause}}%
%BeginExpansion
\pause%
%EndExpansion
\dfrac{15}{-56}=-\dfrac{15}{56}$%
%TCIMACRO{\TeXButton{Pausa}{\pause}}%
%BeginExpansion
\pause%
%EndExpansion

\item $6\cdot \dfrac{5}{11}=\dfrac{6}{1}\cdot \dfrac{5}{11}=\dfrac{6\cdot 5}{%
1\cdot 11}=%
%TCIMACRO{\TeXButton{Pausa}{\pause}}%
%BeginExpansion
\pause%
%EndExpansion
\allowbreak \dfrac{30}{11}$
\end{enumerate}

%TCIMACRO{\TeXButton{EndBlok}{\end{block}}}%
%BeginExpansion
\end{block}%
%EndExpansion

%TCIMACRO{\TeXButton{EndFrame}{\end{frame}}}%
%BeginExpansion
\end{frame}%
%EndExpansion

%TCIMACRO{\TeXButton{BeginFrame}{\begin{frame}}}%
%BeginExpansion
\begin{frame}%
%EndExpansion

\QTR{frametitle}{Conjuntos Num\'{e}ricos}

%TCIMACRO{\TeXButton{BeginBlok}{\begin{block}{Propiedades}}}%
%BeginExpansion
\begin{block}{Propiedades}%
%EndExpansion

\begin{enumerate}
\item Dos fracciones son iguales si, y s\'{o}lo si,el producto de sus
extremos es igual al producto de sus medios, es decir%
\begin{equation*}
\frac{a}{b}=\frac{c}{d}\Longleftrightarrow a\cdot d=b\cdot d
\end{equation*}

\textbf{Ejemplo}. $\dfrac{6}{10}=\dfrac{9}{15},$ ya que $6\cdot 15=10\cdot
9=90$

\item Simplificar una fracci\'{o}n consiste en cancelar un factor com\'{u}n
del numerador y del denominador.

\textbf{Ejemplo.} $\ \dfrac{8}{12}=\dfrac{2\cdot 4}{3\cdot 4}=\dfrac{2}{3}$,
simplicando por el factor com\'{u}n 4.

\item Amplificar una fracci\'{o}n consiste en multiplicar el numerador y
denominador por un mismo n\'{u}mero.

\textbf{Ejemplo.} $\ \dfrac{-3}{8}=\dfrac{-3\cdot 4}{8\cdot 4}=\dfrac{-12}{24%
}$, amplificando por 4.
\end{enumerate}

%TCIMACRO{\TeXButton{EndBlok}{\end{block}}}%
%BeginExpansion
\end{block}%
%EndExpansion

%TCIMACRO{\TeXButton{EndFrame}{\end{frame}}}%
%BeginExpansion
\end{frame}%
%EndExpansion

%TCIMACRO{\TeXButton{BeginFrame}{\begin{frame}}}%
%BeginExpansion
\begin{frame}%
%EndExpansion

\QTR{frametitle}{Conjuntos Num\'{e}ricos}

%TCIMACRO{\TeXButton{BeginBlok}{\begin{block}{Ejercicios}}}%
%BeginExpansion
\begin{block}{Ejercicios}%
%EndExpansion

Multiplique los siguientes n\'{u}meros racionales,\bigskip

\begin{tabular}{lll}
a) $\dfrac{3}{9}\cdot \left( \dfrac{-4}{3}\right) =\allowbreak $%
%TCIMACRO{\TeXButton{Pausa}{\pause}}%
%BeginExpansion
\pause%
%EndExpansion
$-\dfrac{4}{9}$ &  & d) $\dfrac{2}{7}\cdot \dfrac{1}{2}\cdot (-3)=$ 
%TCIMACRO{\TeXButton{Pausa}{\pause}}%
%BeginExpansion
\pause%
%EndExpansion
$-\frac{3}{7}$ \\ 
b) $\left( \dfrac{(-2)-1}{3}\right) \cdot \dfrac{2}{3}=$ 
%TCIMACRO{\TeXButton{Pausa}{\pause}}%
%BeginExpansion
\pause%
%EndExpansion
$-\dfrac{2}{3}$ &  & e) $\dfrac{(-2)\cdot 6}{-3}\cdot \dfrac{(-12)(-9)}{4}=$%
%TCIMACRO{\TeXButton{Pausa}{\pause}}%
%BeginExpansion
\pause%
%EndExpansion
$108$ \\ 
c) $\left( \dfrac{-6}{3}\right) \cdot \left( \dfrac{-4}{7}\right) \cdot 
\dfrac{1}{9}=$%
%TCIMACRO{\TeXButton{Pausa}{\pause}}%
%BeginExpansion
\pause%
%EndExpansion
$\dfrac{8}{63}$ &  & f) $\left( \dfrac{-7}{8}\right) \cdot \left( \dfrac{-4}{%
-3}\right) \cdot \left( \dfrac{2}{14}\right) \cdot (-4)=$ 
%TCIMACRO{\TeXButton{Pausa}{\pause}}%
%BeginExpansion
\pause%
%EndExpansion
$\frac{2}{3}$%
\end{tabular}

%TCIMACRO{\TeXButton{EndBlok}{\end{block}}}%
%BeginExpansion
\end{block}%
%EndExpansion

%TCIMACRO{\TeXButton{EndFrame}{\end{frame}}}%
%BeginExpansion
\end{frame}%
%EndExpansion

%TCIMACRO{\TeXButton{BeginFrame}{\begin{frame}}}%
%BeginExpansion
\begin{frame}%
%EndExpansion

\QTR{frametitle}{Conjuntos Num\'{e}ricos}

%TCIMACRO{\TeXButton{BeginBlok}{\begin{block}{Divisi\'{o}n}}}%
%BeginExpansion
\begin{block}{Divisi\'{o}n}%
%EndExpansion

Para dividir dos n\'{u}meros racionales se multiplica la primera fracci\'{o}%
n por la segunda invertida.

\begin{equation*}
\frac{a}{b}\div \frac{c}{d}=\frac{a}{b}\cdot \frac{d}{c}
\end{equation*}

\textbf{Ejemplos.}

\begin{enumerate}
\item $\dfrac{3}{5}\div \dfrac{2}{9}=\dfrac{3}{5}\cdot \dfrac{9}{2}%
%TCIMACRO{\TeXButton{Pausa}{\pause}}%
%BeginExpansion
\pause%
%EndExpansion
=\allowbreak \dfrac{27}{10}$%
%TCIMACRO{\TeXButton{Pausa}{\pause}}%
%BeginExpansion
\pause%
%EndExpansion

\item $\dfrac{1}{2}\div \left( \dfrac{-2}{3}\right) =\dfrac{1}{2}\cdot
\left( \dfrac{3}{-2}\right) =%
%TCIMACRO{\TeXButton{Pausa}{\pause}}%
%BeginExpansion
\pause%
%EndExpansion
-\dfrac{3}{4}$%
%TCIMACRO{\TeXButton{Pausa}{\pause}}%
%BeginExpansion
\pause%
%EndExpansion

\item $7\div \dfrac{14}{5}=\dfrac{7}{1}\cdot \dfrac{5}{14}=%
%TCIMACRO{\TeXButton{Pausa}{\pause}}%
%BeginExpansion
\pause%
%EndExpansion
\allowbreak \dfrac{5}{2}$%
%TCIMACRO{\TeXButton{Pausa}{\pause}}%
%BeginExpansion
\pause%
%EndExpansion
\end{enumerate}

%TCIMACRO{\TeXButton{EndBlok}{\end{block}}}%
%BeginExpansion
\end{block}%
%EndExpansion

%TCIMACRO{\TeXButton{EndFrame}{\end{frame}}}%
%BeginExpansion
\end{frame}%
%EndExpansion

%TCIMACRO{\TeXButton{BeginFrame}{\begin{frame}}}%
%BeginExpansion
\begin{frame}%
%EndExpansion

\QTR{frametitle}{Conjuntos Num\'{e}ricos}

%TCIMACRO{\TeXButton{BeginBlok}{\begin{block}{Ejercicios}}}%
%BeginExpansion
\begin{block}{Ejercicios}%
%EndExpansion

Divida las siguientes expresiones, simplificando lo m\'{a}s posible

\begin{tabular}{lll}
a) $\dfrac{-8}{5}\div \dfrac{4}{7}=$%
%TCIMACRO{\TeXButton{Pausa}{\pause} }%
%BeginExpansion
\pause
%EndExpansion
&  & d) $\dfrac{7}{9}\div \dfrac{-4}{12}=%
%TCIMACRO{\TeXButton{Pausa}{\pause}}%
%BeginExpansion
\pause%
%EndExpansion
\allowbreak $ \\ 
b) $\left( \dfrac{-1}{3}\right) \div \left( \dfrac{2}{-3}\right) =$%
%TCIMACRO{\TeXButton{Pausa}{\pause} }%
%BeginExpansion
\pause
%EndExpansion
&  & e) $\left( \dfrac{(-2)\cdot 6}{-3}\right) \div \left( \dfrac{(-15)(-2)}{%
4}\right) =$%
%TCIMACRO{\TeXButton{Pausa}{\pause} }%
%BeginExpansion
\pause
%EndExpansion
\\ 
c) $\left( \dfrac{-3}{4}\right) \div 6=\allowbreak $%
%TCIMACRO{\TeXButton{Pausa}{\pause} }%
%BeginExpansion
\pause
%EndExpansion
&  & f) $\left( \left( \dfrac{-5}{8}\right) \cdot \left( \dfrac{-4}{-5}%
\right) \right) \div \left( \dfrac{12}{6}\right) =$%
\end{tabular}

%TCIMACRO{\TeXButton{EndBlok}{\end{block}}}%
%BeginExpansion
\end{block}%
%EndExpansion

%TCIMACRO{\TeXButton{EndFrame}{\end{frame}}}%
%BeginExpansion
\end{frame}%
%EndExpansion

%TCIMACRO{\TeXButton{BeginFrame}{\begin{frame}}}%
%BeginExpansion
\begin{frame}%
%EndExpansion

\QTR{frametitle}{Ejercicios}

%TCIMACRO{\TeXButton{BeginBlok}{\begin{block}{Ejemplos}}}%
%BeginExpansion
\begin{block}{Ejemplos}%
%EndExpansion

Calcular las siguientes expresiones:

\begin{multicols}{2}

\begin{enumerate}

\item $(-5+(-7))=\allowbreak $

\item $(3+(-4)-8+12)=\allowbreak \allowbreak $

\item $\dfrac{12(-40)(-12)}{-5(-3)-3(-3)}=\allowbreak $

\item $\dfrac{(-3)(8)(-2)}{(-4)(-6)-2(-12)}=\allowbreak $

\item $\dfrac{-2}{9}+\dfrac{4}{7}-\dfrac{1}{2}=\allowbreak $

\item $\dfrac{\frac{2}{3}+\frac{1}{2}}{\frac{2}{2}}=\allowbreak \allowbreak $

\item $\dfrac{\frac{1}{8}+\dfrac{1-\dfrac{1}{2}}{2}}{\frac{1}{12}}%
=\allowbreak $

\item $4+\dfrac{3}{5}-\dfrac{7}{2}-\dfrac{4}{12}-\dfrac{1}{6}=\allowbreak $

\item $\dfrac{(-5)\left( \frac{4}{7}\right) }{(-2)\left( \frac{-3}{5}\right) 
}=\allowbreak $

\end{enumerate}

\end{multicols}

%TCIMACRO{\TeXButton{EndBlok}{\end{block}}}%
%BeginExpansion
\end{block}%
%EndExpansion

%TCIMACRO{\TeXButton{EndFrame}{\end{frame}}}%
%BeginExpansion
\end{frame}%
%EndExpansion

\end{document}
